% Fixes long, unbreakable citation URLs overflowing the page margin in exported PDFs.
%
% Pandoc's own default LaTeX template already tries to solve exactly this problem -- its
% `after-header-includes.latex` partial contains:
%
%   \IfFileExists{xurl.sty}{\usepackage{xurl}}{} % add URL line breaks if available
%
% -- but `xurl.sty` is not part of this app's bundled TinyTeX distribution (confirmed: zero
% matches for `xurl*` anywhere under Resources/TinyTeX), so that conditional silently no-ops.
% Pandoc falls back to plain `url.sty`'s much narrower default break points (mostly
% punctuation like / . - ~ : ? #), which don't cover the long, unbroken alphanumeric runs
% common in DOI/tracking-parameter URLs -- exactly what a reported PDF export showed
% overflowing past the page margin (confirmed: without this fix, a long unbroken-alphanumeric
% URL produces a LaTeX "Overfull hbox" hundreds of points wide; with it, well under 3pt --
% normal justification slack).
%
% Rather than vendor xurl.sty itself into TinyTeX -- this project has no package manifest or
% update script for TinyTeX (see docs/architecture/export.md's "Bundled TinyTeX" section), and
% the one other missing-package case found in this codebase (`caption`, see the comment atop
% float-package.tex) was likewise solved with a small header workaround instead of bundling a
% new package -- this reproduces xurl's documented effect directly via `url.sty`'s own public
% `\UrlBreaks` extension point. Per CTAN's package description (https://ctan.org/pkg/xurl):
% "This package loads url by default and defines possible URL breaks for all alphanumerical
% characters, as well as = / . : * - ~ ' \"".
%
% `\usepackage{url}` here is deliberate even though pandoc's template also loads url.sty
% indirectly (via `\usepackage{bookmark}` requiring hyperref, which requires url) later in the
% preamble, in `after-header-includes.latex` -- `--include-in-header` content is inserted
% earlier than that partial, so `\UrlBreaks` would not exist yet to extend below if we
% depended on that later load. Loading `url` twice with compatible options is a standard,
% harmless LaTeX no-op (packages guard against double-loading).
%
% `\g@addto@macro` is a LaTeX2e kernel macro (not from url.sty or any extra package) used here
% to *append* to the existing `\UrlBreaks` list rather than replace it, so url.sty's own
% default punctuation break points are preserved alongside the new ones. The `\do\<char>`
% list form matches the LaTeX kernel's own `\dospecials`-style idiom for registering
% characters (including letters) as catcode-independent tokens -- this is the standard,
% documented way to extend `\UrlBreaks`.
\usepackage{url}
\makeatletter
\g@addto@macro{\UrlBreaks}{%
  \do\a\do\b\do\c\do\d\do\e\do\f\do\g\do\h\do\i\do\j\do\k\do\l\do\m%
  \do\n\do\o\do\p\do\q\do\r\do\s\do\t\do\u\do\v\do\w\do\x\do\y\do\z%
  \do\A\do\B\do\C\do\D\do\E\do\F\do\G\do\H\do\I\do\J\do\K\do\L\do\M%
  \do\N\do\O\do\P\do\Q\do\R\do\S\do\T\do\U\do\V\do\W\do\X\do\Y\do\Z%
  \do\0\do\1\do\2\do\3\do\4\do\5\do\6\do\7\do\8\do\9%
  \do\=\do\:\do\*\do\~\do\'\do\"%
}
\makeatother
